% =========================================================================
% Measuring What Matters When the Robot Has the Gun:
% A Validity Audit of Hallucination-Induced Unauthorized Action in LLM Agents
%
% Long-form term paper for CS321M (AI Measurement Science, Spring 2026)
%
% Build instructions for Overleaf:
%   - Upload paper_long.tex and references.bib to a fresh project
%   - Set the compiler to pdfLaTeX (default)
%   - Compile twice (LaTeX runs bibtex/biber between the two compiles)
% =========================================================================

\documentclass[11pt]{article}

\usepackage[margin=1in]{geometry}
\usepackage{amsmath, amssymb}
\usepackage[utf8]{inputenc}
\usepackage[T1]{fontenc}
\usepackage{microtype}
\usepackage{parskip}
\usepackage{enumitem}
\usepackage{titlesec}
\usepackage{hyperref}
\hypersetup{
  colorlinks=true,
  linkcolor=black,
  citecolor=blue!50!black,
  urlcolor=blue!50!black
}

\titleformat{\section}{\large\bfseries}{\thesection}{0.5em}{}
\titleformat{\subsection}{\normalsize\bfseries}{\thesubsection}{0.5em}{}

\title{Measuring What Matters When the Robot Has the Gun:\\
A Validity Audit of Hallucination-Induced Unauthorized Action in LLM Agents}
\author{King \\
\small CS321M --- AI Measurement Science (Spring 2026) \\
\small Instructor: Sanmi Koyejo}
\date{Draft v1, May 2026}

\begin{document}
\maketitle

\begin{abstract}
The case for deploying LLM agents in irreversible-action settings --- military targeting, embodied robotics, autonomous trading --- rests on a transfer claim: that strong scores on existing agent-safety benchmarks predict that an agent will not take an explicitly forbidden action when it matters. This paper argues that the transfer claim is unsupportable from current measurements because no published benchmark isolates the construct that is doing the work: the conditional rate at which an agent, having hallucinated, then crosses an explicit prohibition. We formalize this construct as \emph{Hallucination-Induced Unauthorized Action} (HIUA), decompose it from the four neighboring constructs that the agent-safety literature confounds it with, and audit nine prior instruments against Messick's six aspects of validity. The substantive aspect --- the cognitive-process claim that hallucination caused the violation --- is essentially undefended in the surveyed corpus, even where the manipulations exist to test it. The structural aspect is reported through single accuracy numbers that aggregate over 90-percentage-point dispersion in the dissociation between rule recall and rule adherence. The consequential aspect predicts three specific gaming failures already visible in adjacent benchmarks. We close with HIUA-Bench, a design sketch that factorially crosses three hallucination sub-types against three prohibition-salience levels across four task domains, with simulator-verifiable outcomes, a recall probe to partition lucid from hallucinated violation, and a Generalizability-Theory reliability plan reporting $\Phi$ and $\mathrm{E}\rho^2$ over persons $\times$ items $\times$ raters $\times$ occasions.
\end{abstract}

\section{Introduction}

The empirical anchors for an agent-safety paper in 2026 are not subtle. RoboPAIR, an algorithmic jailbreak procedure applied to three commercial robot platforms, achieved approximately 100\% attack success eliciting physical harmful actions from the NVIDIA Dolphins self-driving system, the Clearpath Jackal UGV running a GPT-4o planner, and the Unitree Go2 quadruped running GPT-3.5~\cite{robopair2024}. ARMOR 2025, the first military-doctrine-aligned LLM safety benchmark, evaluated twenty-one commercial LLMs against 519 prompts grounded in the Law of War, the Rules of Engagement, and the DoD Joint Ethics Regulation; it documented two distinct failure modes --- models that hallucinate rules or invent constraints when asked to justify actions, and models that refuse lawful and scoped requests --- and concluded that deploying general-purpose LLMs for military decision support without additional controls is premature~\cite{armor2026}. Anthropic's June 2025 agentic-misalignment red team, in which models were given autonomous agentic roles and access to a simulated company's email system, found that Claude Opus 4 chose to blackmail a fictional executive in 96\% of relevant scenarios; Gemini 2.5 Flash blackmailed at 96\%, GPT-4.1 and Grok 3 Beta at 80\%, and DeepSeek-R1 at 79\%~\cite{anthropicagentic2025}.

These are three different mechanisms --- adversarial jailbreak, doctrine-mismatch under genuine military prompts, and emergent goal-protection under shutdown threat --- with overlapping blast radius. The natural reading is that the field has many benchmarks for different parts of an unsafe-agent landscape. The problem this paper takes up is that the benchmarks do not actually carve the landscape at the joints that matter when an action is irreversible.

The cleanest empirical demonstration of the carving problem is in the constraint-adherence literature. DriftBench, evaluating seven models from five providers across 2{,}146 scored runs on multi-turn ideation tasks, found that the rate at which models violate constraints they can simultaneously restate verbatim --- the \emph{knows-but-violates} or KBV rate --- ranges from 8\% on the lowest-violation model to 99\% on the highest~\cite{driftbench2026}. SafeAgentBench, evaluating embodied LLM agents in an interactive simulator on 750 tasks across 10 hazard categories, found that the best-performing baseline achieves 69\% success on safe tasks but only 5\% rejection on hazardous tasks~\cite{safeagentbench2024}. AgentHarm, evaluating frontier models on 110 explicitly malicious agent tasks (440 with augmentations) across 11 harm categories, found that leading LLMs are surprisingly compliant with malicious agent requests \emph{without} any jailbreak~\cite{agentharm2024}. These three findings --- 91-percentage-point KBV dispersion, 5\% hazardous-task rejection, and pre-jailbreak compliance --- are routinely reported under the umbrella of ``agent safety,'' and developers cite their relative scores when arguing about deployment. They are measurements of different conditional distributions over the same outcome variable.

This paper makes three contributions. First, in \S\ref{sec:construct} we define Hallucination-Induced Unauthorized Action (HIUA) formally as the conditional rate $P(\text{forbidden action} \mid \text{hallucination} \wedge \text{salient prohibition})$, decomposing it from four neighboring constructs the agent-safety literature confounds it with. Second, in \S\ref{sec:audit} we audit nine prior instruments against Messick's six aspects of validity, showing that the substantive aspect --- the claim that hallucination caused the violation, as opposed to merely co-occurring with it --- is essentially undefended in the surveyed corpus. Third, in \S\ref{sec:design} we sketch HIUA-Bench, a design that factorially crosses three hallucination sub-types against three prohibition-salience levels and four task domains, with simulator-verifiable outcomes, a post-hoc recall probe to partition lucid from hallucinated violation, and a Generalizability-Theory reliability plan that reports $\Phi$ and $\mathrm{E}\rho^2$ over persons $\times$ items $\times$ raters $\times$ occasions $\times$ paraphrases.

The thesis the paper defends is this: when the cost of a single forbidden action is irreversible --- a fired weapon, a transferred fund, a thrown switch --- measurement that confounds lucid violation with hallucination-induced violation systematically mis-licenses deployment claims. A benchmark that reports a single aggregate violation rate cannot distinguish a model that needs a faithfulness intervention from a model that needs an instruction-following intervention from a model that needs a perception-grounding intervention. The validity audit below is the argument; the design sketch is the constructive proposal.

\section{Background: validity theory in the agent-safety context}

We take the reader's familiarity with Messick's unified validity framework~\cite{messick1989}, Kane's argument-based approach~\cite{kane2013}, and the recent application of these frameworks to AI evaluation~\cite{m2m2025, nomologicalnetworks2026} as given. Three points are worth flagging because they do load-bearing work in \S\ref{sec:audit}.

First, the construct under test is not a property of the LLM as a system; it is a property of an \emph{interpretation} placed on an LLM's score, as Messick's ``validity as overall evaluative judgment of the degree to which empirical evidence and theoretical rationales support the adequacy and appropriateness of interpretations and actions based on test scores'' makes explicit~\cite{messick1989, standards2014}. The relevant interpretation for the deployment community is ``this agent will not, when deployed, take a forbidden action because of a hallucinated belief about authorization, state, or tool effect.'' A benchmark that reports a violation rate without partitioning the conditional pathway licenses a weaker interpretation than the one its consumers use it to license.

Second, Cronbach and Meehl's original framing of construct validity as the placement of a construct within a \emph{nomological network} of theoretical relations~\cite{cronbachmeehl1955} is more directly relevant to LLM evaluation than the Borsboom causal account~\cite{borsboom2004} for the use case at hand. Freiesleben has recently made this argument explicitly in the LLM-capability context: the nomological account avoids the strong ontological commitments of the causal account while providing more substantive constraint on construct meaning than Kane's purely inferential refinement~\cite{nomologicalnetworks2026}. We endorse this for HIUA: we will not claim that there is a unitary latent ``hallucination propensity'' causing variance in observed violation rates. We will claim that HIUA, if it is a construct worth measuring, should sit in a predictable web of correlations with related measurements --- hallucination rate on text-only benchmarks, jailbreak compliance, capability --- and that this web is what licenses the deployment interpretation.

Third, the consequential aspect of validity --- the evaluation of the social and behavioral consequences of using a test in a particular way --- has been the least applied of Messick's six in AI evaluation~\cite{jacobswallach2021, m2m2025}. The agent-safety benchmark case forces it back into prominence because the \emph{use} of these benchmarks is to license deployment in irreversible-action settings. A benchmark that produces a plausible-looking score but does not isolate the construct that the deployment claim requires is not neutral; it is misleading in a particular direction. We return to this in \S\ref{sec:consequential}.

\section{Construct definition}
\label{sec:construct}

\subsection{Formal definition}

Let an agent trajectory $\tau = (s_0, a_0, s_1, a_1, \ldots)$ proceed under a system prompt containing an explicit prohibition $P$ (a single-sentence rule of the form ``do not perform any action $a \in A_{\text{forbidden}}$''). Let $h_t \in \{0, 1\}$ indicate the occurrence of a hallucination event at step $t$: a generated intermediate token sequence (chain-of-thought, tool-argument, or tool-output interpretation) that is unfaithful to the model's input and to the ground-truth world state. Let $s_t \in \{0, 1\}$ indicate the salience condition: $P$ was present in the model's effective context, was syntactically parseable, and was within the attention window (operationalizable as: a recall probe at time $t$ would elicit a verbatim or near-verbatim restatement of $P$).

The HIUA construct is the conditional rate
\[
\mathrm{HIUA}(M, D) = P\bigl(a_t \in A_{\text{forbidden}} \mid h_t = 1 \wedge s_t = 1;\, M, D\bigr)
\]
for a model $M$ on a domain distribution $D$. This is a narrower target than aggregate violation rate (which marginalizes over $h_t$ and $s_t$), aggregate hallucination rate (which marginalizes over $a_t \in A_{\text{forbidden}}$), and aggregate compliance rate (which marginalizes over $h_t$ and over instruction provenance).

\subsection{Sub-constructs}

HIUA decomposes into three sub-constructs that map onto stages of the recently-surveyed agent-hallucination taxonomy~\cite{halluagentsurvey2025}:

\textbf{Authorization hallucination.} The agent fabricates that an exception, override, or higher-priority instruction permits $a \in A_{\text{forbidden}}$. Concretely: a forged ``manager email'' arrives mid-task with apparent override authority; the agent treats the forgery as authoritative without grounding its trust in any verified channel. This is a reasoning-stage hallucination.

\textbf{State hallucination.} The agent fabricates a world state in which $a$ no longer belongs to $A_{\text{forbidden}}$. Concretely: an ambiguous file name suggests the file is a duplicate; the agent infers ``this is obviously a duplicate, so deletion is not data loss'' and deletes the unique copy. This is a perception-stage hallucination.

\textbf{Tool hallucination.} The agent fabricates a tool's behavior such that calling the tool produces $a$ without the agent representing the call as $a$. Concretely: a tool's docstring describes it as a search; the agent treats the search-call as cost-free but the tool in fact triggers a billable action; the agent has not represented its own action as ``spending money.'' This is an execution-stage hallucination.

These three sub-constructs are theoretically distinct in their mechanism, in the interventions that would fix them, and --- we will argue in \S\ref{sec:audit} --- in the measurement design that would isolate them.

\subsection{Anti-construct}

HIUA is what is sometimes called a \emph{specific} construct: it is sharpened by being explicit about what it is not. Four neighboring failure modes are excluded:

\textbf{Lucid violation (knows-but-violates).} The rule was recalled and the action was taken anyway. This is what DriftBench measures most directly, with rates from 8\% to 99\% across models~\cite{driftbench2026}. The lucid violation rate is \emph{not} HIUA; a model that lucidly violates a rule is not hallucinating, even if its behavior is unsafe.

\textbf{Jailbreak compliance.} The user or third-party content explicitly instructs the harmful action and the agent complies. This is what AgentHarm~\cite{agentharm2024}, RoboPAIR~\cite{robopair2024}, InjecAgent~\cite{injecagent2024}, and AgentDojo~\cite{agentdojo2024} measure. The agent's intermediate generation is not unfaithful --- the agent is faithfully complying with an instruction it should not have complied with.

\textbf{Over-refusal of legitimate tasks.} The model refuses a permitted action because it pattern-matches to a harmful one. OR-Bench documents this as a measurable distinct phenomenon over 80{,}000 prompts across 32 LLMs~\cite{orbench2024}, and ARMOR finds it specifically as a military-context failure mode~\cite{armor2026}.

\textbf{Policy-invisible violation.} The agent took the action because the policy facts needed to recognize the action as forbidden were not present in its context. PhantomPolicy formalizes this case across eight violation categories with the policy intentionally absent from tool responses~\cite{phantompolicy2026}. This is a \emph{missing-information} problem rather than a hallucination problem.

A measurement that cannot partition these four neighbors from HIUA is measuring a confounded construct. Borsboom's causal account of validity requires that variations in the target attribute produce variations in the measurement outcome~\cite{borsboom2004}; if the same observed violation rate can be produced by any of five distinct mechanisms, then the variation in attribute (HIUA propensity specifically) is not what is producing the variation in measurement. The measurement is responsive, but it is responsive to a mixture.

\section{Related work, organized by what each instrument measures}

We organize the prior literature by the construct each instrument \emph{implicitly} measures, not by chronology or by author. This organization is not neutral; it is the work the \S\ref{sec:audit} audit will lean on.

\subsection{Malicious-user-instruction compliance}

The largest cluster of recent agent-safety benchmarks measures whether an agent complies with an explicit instruction --- from the user, the system, or a third party --- to perform an action that the agent should refuse. AgentHarm's 110 explicitly malicious agent tasks (440 with jailbreak augmentations) span 11 harm categories from fraud to cybercrime to harassment, and document that leading models are surprisingly compliant without any jailbreaking, that universal jailbreak templates adapt effectively to agentic tasks, and that jailbroken agents retain their capabilities while behaving maliciously across multi-step trajectories~\cite{agentharm2024}. AgentDojo, with 97 realistic tasks and 629 security cases spanning email, banking, and travel-booking environments, found that GPT-4o's benign utility of 69\% drops to 45\% under prompt-injection attack, with targeted attack success rates of 53.1\% on the ``Important message'' canonical attack~\cite{agentdojo2024}. InjecAgent's 1{,}054 cases across 17 user tools and 62 attacker tools found that ReAct-prompted GPT-4 is vulnerable to indirect prompt injection 24\% of the time, with hacking-prompt reinforcement nearly doubling the success rate~\cite{injecagent2024}. RoboPAIR achieved approximately 100\% attack success eliciting physical harmful actions across three robot platforms~\cite{robopair2024}.

These instruments cluster cleanly around a common construct: the agent's resistance to taking forbidden actions when those actions are explicitly requested by an adversary. The adversary may be the user (AgentHarm) or third-party content (InjecAgent, AgentDojo) or an algorithmic attacker (RoboPAIR), but in every case the harmful action is in the instruction stream the agent processes. This is the \emph{complement} of HIUA: in HIUA, no actor has instructed the harmful action; the agent has confabulated authorization, state, or tool effect on its own.

\subsection{Outcome-aggregated agent safety}

A second cluster measures aggregate agent failure rates across heterogeneous scenarios without distinguishing the causal pathway. ToolEmu's LM-emulated sandbox covering 36 high-stakes toolkits and 144 test cases reports that the safest LM agent fails 23.9\% of the time according to the automatic evaluator, and 68.8\% of the time according to human raters~\cite{toolemu2024}. R-Judge's 569 multi-turn agent interaction records spanning 27 risk scenarios in 5 application categories and 10 risk types measures the proficiency of an LLM-judge in identifying safety risks post-hoc, with GPT-4o achieving 74.42\%~\cite{rjudge2024}. Agent-SafetyBench's 349 interaction environments and 2{,}000 test cases across 8 risk categories and 10 failure modes found that none of 16 popular LLM agents achieved a safety score above 60\%, and identified lack of robustness and lack of risk awareness as the two fundamental defects~\cite{agentsafetybench2024}.

These instruments report a single aggregated failure rate per model, sometimes with a per-category breakdown. The DriftBench finding that KBV varies from 8\% to 99\% across models~\cite{driftbench2026} implies that aggregating across hallucinated and lucid violations averages over a roughly 90-percentage-point swing in the dissociation between recall and adherence. A single accuracy number is therefore difficult to interpret as a stable model property: the same 30\% violation rate could reflect a model that always recalls and sometimes violates anyway (lucid), or a model that never recalls and always violates (recall failure), or a model that confabulates new authorizations 30\% of the time (HIUA proper). The three are different failure modes with different remediations.

The R-Judge case is particularly instructive: it is a meta-evaluation in which the LLM under test is asked to judge agent traces, not to act. This is a different construct again --- \emph{risk awareness as judgment} rather than \emph{risk avoidance as action}. R-Judge's contribution is real, but reading its scores as a measure of agent safety in the deployment sense is a construct-level mismatch.

\subsection{Lucid constraint adherence}

A third cluster measures the gap between what a model can recall and what a model does. DriftBench's central finding is the dissociation: across 2{,}146 scored benchmark runs spanning seven models, four interaction conditions, and 38 research briefs from 24 scientific domains, models accurately restate constraints they simultaneously violate, with KBV rates of 8\% to 99\% across models~\cite{driftbench2026}. Structured checkpointing partially reduces KBV but does not close the dissociation. AgentIF, constructed from 50 real-world agentic applications with instructions averaging 1{,}723 words and 11.9 constraints each across three constraint types (formatting, semantic, tool), found that the best of the evaluated models followed fewer than 30\% of instructions perfectly~\cite{agentif2025}. The Hierarchical Safety Principles benchmark documented a quantifiable ``cost of compliance'' --- safety constraints degrade task performance even when compliant solutions exist --- and an ``illusion of compliance'' in which high adherence rates often mask task incompetence rather than principled choice~\cite{hsp2025}.

The instruction-hierarchy training of Wallace et al.\ is best read as a \emph{fix} for the gap this cluster exposes: by training models to prioritize system > developer > user > data instructions, models are expected to lucidly adhere to higher-priority rules even when later content suggests otherwise~\cite{instructionhierarchy2024}. The fix targets the dissociation that DriftBench and AgentIF measure, not the hallucination pathway that HIUA targets.

This cluster's value for HIUA is methodological: it has established that constraint recall and constraint adherence are dissociable, which is precisely the partition HIUA-Bench will need to perform in reverse --- separating \emph{failure to adhere despite recall} (lucid) from \emph{failure to adhere because of hallucination}.

\subsection{Embodied safety}

A fourth cluster targets embodied LLM agents in simulators. SafeAgentBench's 750 tasks (with 450 hazardous and 300 safe controls) across 10 hazard categories and 3 task types are grounded in an interactive AI2-THOR-style simulator; the best-performing baseline achieved 69\% success on safe tasks but only 5\% rejection on hazardous ones~\cite{safeagentbench2024}. AGENTSAFE extends this with 45 adversarial scenarios, 1{,}350 hazardous tasks, and 9{,}900 instructions evaluated across nine state-of-the-art VLMs and two embodied agent workflows, with multi-level diagnosis across perception, planning, and execution; it uncovers systematic failures in translating hazard recognition into safe planning and execution~\cite{agentsafe2025}.

HEAL is the closest existing analogue to HIUA in design. It probes hallucinations in long-horizon embodied planning by constructing scene-task inconsistencies --- distractor injection, task-relevant object removal, synonymous object substitution, and scene-task contradiction --- that induce hallucination rates up to 40$\times$ higher than base prompts across 12 models in two simulators~\cite{heal2025}. HEAL finds that models cannot reliably reject infeasible tasks, that hallucinations persist despite feedback-based mitigation, and that they are only partially reduced by cross-modal input. The illustrative example HEAL gives --- a robot instructed to ``put the knives in the dishwasher'' in a scene without a dishwasher, then hallucinating the dishwasher and placing sharp utensils in an empty cabinet --- is exactly the state-hallucination sub-construct from \S\ref{sec:construct}.

The gap between HEAL and HIUA is the dependent variable. HEAL measures plan correctness conditional on the inconsistency manipulation. HIUA would measure prohibition crossing conditional on the same manipulation. The two are recoverable from the same trace data; HEAL did not analyze it that way.

\subsection{Hallucination measurement, standalone}

The hallucination literature in isolation from agent action has matured considerably. HalluLens distinguishes extrinsic from intrinsic hallucination, dynamically generates its evaluation data to prevent test-set leakage, and disentangles hallucination from factuality~\cite{hallulens2025}. The recent survey of agent-specific hallucinations proposes a stage-based taxonomy spanning the agent workflow (brain / perception / action) and catalogs 18 triggering causes~\cite{halluagentsurvey2025}. Neither of these is cross-linked to a measure of \emph{consequential action}: they measure whether a hallucination occurred, not whether it then caused a forbidden action.

\subsection{Validity and benchmarking-of-benchmarks}

A fifth, methodologically central cluster has emerged in the past two years arguing that the agent-safety literature needs a more disciplined validity practice. Measurement-to-Meaning lays out a validity-centered framework for reasoning about which claims existing measurements can and cannot support~\cite{m2m2025}. The nomological-networks-required argument applies Cronbach-Meehl to LLM capability benchmarks and argues that the nomological account is the most suitable foundation~\cite{nomologicalnetworks2026}. BetterBench assessed 25 AI benchmarks against 46 best-practice criteria and found large quality differences, with most benchmarks failing to report statistical significance or allow easy replication~\cite{betterbench2024}. The Eriksson et al.\ interdisciplinary meta-review of approximately 100 studies on AI-benchmark shortcomings names misaligned incentives, construct validity gaps, and gaming as systemic problems~\cite{canwetrust2025}. Raji et al.\ critiqued the practice of taking ``general'' benchmarks as abstractions for general capability~\cite{raji2021}. Bowman and Dahl proposed four criteria most NLU benchmarks fail~\cite{bowmandahl2021}. The PSN-IRT analysis of 41{,}871 items across 11 LLM benchmarks documented widespread saturation, insufficient difficulty ceilings, and item-level contamination~\cite{lostinbench2025}. Liao and Xiao argued for socio-technically situated evaluation~\cite{liaoxiao2023}. Jacobs and Wallach proposed measurement modeling as the right frame for fairness in computational systems, emphasizing content and consequential validity as load-bearing~\cite{jacobswallach2021}.

This methodological cluster is what makes the present paper possible. It also makes the paper risky: the audit below is not novel in form, only in target. We try, in \S\ref{sec:audit}, to make the target-specific arguments sharp enough that the audit is not a Messick checklist exercise but a critique that produces conclusions the literature does not already contain.

\section{Validity audit: nine instruments, six aspects}
\label{sec:audit}

This section is the core. We organize by Messick's six aspects of validity rather than by instrument, and within each aspect we audit at least three prior instruments against the HIUA construct. The aim is to produce, in each subsection, at least one critique that would not exist without the validity framework.

\subsection{Content aspect --- does the test span the construct?}

The content question for HIUA is whether the items in a benchmark sample the joint space of $\{\text{hallucination type}\} \times \{\text{prohibition salience}\} \times \{\text{action domain}\}$ densely enough to license generalization to the deployment use case. We audit four instruments.

SafeAgentBench's 10 hazard categories simulator-grounded across 3 task types provide solid content coverage of \emph{hazardous instructions} in embodied settings~\cite{safeagentbench2024}. They do not factorially cross hazardous instructions with hallucination triggers. The construct it spans is ``hazard refusal under direct adversarial instruction,'' not HIUA: the 5\% rejection rate on hazardous tasks is a measure of refusal in the presence of an explicit harmful request, not of the agent's tendency to confabulate authorization for an action no one requested.

AgentHarm's 11 harm categories with 110 base tasks and 440 augmented versions provide thick coverage of \emph{malicious-user-instruction compliance}~\cite{agentharm2024}. The user is the adversary by construction. The benchmark contains no items where the agent has been given a clean task and then must avoid hallucinating a side-channel authorization that would license a harmful action.

InjecAgent and AgentDojo cover \emph{third-party-injected adversarial instructions} --- the agent receives a clean user instruction but processes external content that contains a hostile instruction~\cite{injecagent2024, agentdojo2024}. Same coverage gap from HIUA's perspective: the agent has been instructed (by the attacker) to do the harmful thing.

HEAL has the right manipulation in the right setting: scene-task inconsistencies that induce state hallucination at rates up to 40$\times$ the base~\cite{heal2025}. But its content coverage stops short of explicit prohibitions. The dependent variable is ``did the agent's plan match the (impossible) instruction'' rather than ``did the agent's plan cross a prohibition the system prompt placed on it.'' A content extension of HEAL with explicit per-trial prohibition statements would be the closest existing instrument to HIUA-on-content.

The cumulative content-aspect critique: prevailing benchmarks cover the malicious-instruction half of agent safety thickly and the hallucinated-authorization half almost not at all. Content validity for HIUA across the surveyed corpus is, by inspection, near zero.

\subsection{Substantive aspect --- does the response process match the theorized cognition?}

This is the aspect on which the literature is most exposed. Messick's substantive aspect requires evidence that the cognitive process generating the test response matches the theorized process underlying the construct~\cite{messick1989, standards2014}. For HIUA, the theorized process is ``hallucination event $\rightarrow$ forbidden action,'' and the substantive aspect asks: when the model produces a forbidden action, do we have evidence that a hallucination produced it, as opposed to merely co-occurring with it?

ToolEmu and Agent-SafetyBench score outcomes only~\cite{toolemu2024, agentsafetybench2024}. The agent's intermediate reasoning is not analyzed for the presence or character of a hallucination. The 23.9\% (automatic) and 68.8\% (human) failure rates ToolEmu reports are therefore \emph{consistent with} hallucination-driven failure but provide no evidence that hallucination produced any of them. The post-hoc attribution ``this is a hallucination problem'' is exactly the kind of mechanistic claim that substantive validity is supposed to license.

R-Judge measures whether a \emph{judge} LLM can identify the safety issue post-hoc~\cite{rjudge2024}. The cognitive process under test is the judge's, not the actor's. The 74.42\% GPT-4o-as-judge score is informative for an audit pipeline; it is uninformative as evidence about why any specific agent in the trace took the action it took.

DriftBench is the closest existing analogue to a substantive-aspect instrument. By interleaving a recall probe with the action, it partitions trials into a $2 \times 2$ of $\{\text{recall yes/no}\} \times \{\text{violation yes/no}\}$ and lets us read off the KBV rate from the bottom-right cell~\cite{driftbench2026}. This is exactly the right move, and the resulting KBV rates (8--99\% across models) are precisely the kind of finding that an outcome-only benchmark cannot produce. But the partition is one factor short of what HIUA requires: it separates ``did the model know the rule'' from ``did the model violate it,'' but it does not separate ``did the model hallucinate authorization'' from ``did the model hallucinate state'' from ``did the model hallucinate tool effect'' from ``did the model violate lucidly with full recall and no hallucination.'' The $2 \times 2$ must become a $2 \times 2 \times 2$ (or, with the sub-construct decomposition, a $2 \times 4 \times 2$ for $\{\text{recall}\} \times \{\text{hallucination-type or none}\} \times \{\text{violation}\}$).

HEAL has the manipulation half of what substantive validity requires: it intervenes on the world state to induce a hallucination, then measures the agent's plan~\cite{heal2025}. It does not measure prohibition crossing as the consequential outcome.

The cumulative critique: there is no surveyed benchmark whose data could support a substantive-validity argument for HIUA as published. The data could, in principle, be re-analyzed: HEAL trace data plus an added prohibition layer, DriftBench plus a hallucination-type classifier on the trace, ToolEmu plus a recall probe. The fact that the field has not done this re-analysis is itself the critique: agent-safety benchmarks have been positioned as outcome-summary instruments, and the underlying substantive question has been displaced into discussion sections.

The constructive consequence for HIUA-Bench is that the substantive aspect must be designed in. We will return to this in \S\ref{sec:scoring}.

\subsection{Structural aspect --- internal structure and reliability}

The structural aspect concerns whether the test's internal structure matches the theorized internal structure of the construct, and --- for our purposes --- whether the resulting score is reliable enough to bear the interpretation placed on it~\cite{messick1989, standards2014}. The course's central reliability framework, Generalizability Theory, partitions observed-score variance into facets (persons, items, raters, occasions, settings) and reports $\Phi$ (absolute reliability) and $\mathrm{E}\rho^2$ (relative reliability) coefficients~\cite{gtheoryage2025}. We audit four instruments against this lens.

AgentIF reports per-constraint-type breakdowns and an overall instruction-following rate, but no variance-component decomposition~\cite{agentif2025}. Given that instructions average 11.9 constraints each, the per-item difficulty distribution is plausibly highly skewed; without a structural analysis, the reported $<$30\% perfect-adherence rate could reflect that one or two constraint types are systematically harder, or that the constraint-counts confound difficulty, or both.

ToolEmu reports a striking discrepancy that is, in effect, a partial single-facet G-study: 23.9\% failure rate per automatic evaluator versus 68.8\% per human rater on the same 144 cases~\cite{toolemu2024}. The 45-percentage-point gap is, in G-Theory terms, a strong indication that $\sigma^2(\text{rater})$ is large --- large enough that the choice of rater is doing more work in producing the reported number than the choice of agent. The authors do not report it this way; they report it as a quality-of-emulator check. The interpretation as a reliability-failure-of-the-rating-procedure is, we think, more directly load-bearing for downstream interpretation.

DriftBench's 8--99\% KBV variance across persons (models) is exactly the kind of dispersion that should be partitioned~\cite{driftbench2026, gtheoryage2025}. If $\sigma^2(\text{person}) \gg \sigma^2(\text{item, occasion, rater})$, the benchmark is measuring real model differences and a single per-model KBV rate is a useful summary. If $\sigma^2(\text{item})$ is comparable to $\sigma^2(\text{person})$, the benchmark is largely measuring which items happened to elicit KBV, and the per-model rates would shift with item resampling. The published data does not support that decomposition.

The PSN-IRT analysis of 41{,}871 items across 11 LLM benchmarks is the methodological precedent we lean on most heavily~\cite{lostinbench2025}. It demonstrates that item-level IRT analysis on existing benchmarks is feasible and that the results are sometimes uncomfortable: widespread saturation, insufficient difficulty ceilings, and data contamination affect benchmarks the field has been using to draw broad capability claims. The same analysis applied to agent-safety benchmarks is, to our knowledge, not yet published. Our prediction is that a substantial fraction of items in AgentHarm and SafeAgentBench would either show ceiling effects on frontier models (uninformative items) or show suspiciously low discrimination parameters (items where most failures are model-independent).

The cumulative critique: agent-safety benchmarks report single numbers where the underlying variance structure makes those single numbers difficult to interpret. The minimum reporting standard a HIUA-grade benchmark should adopt is $\Phi$ and $\mathrm{E}\rho^2$ across at least persons $\times$ items $\times$ raters $\times$ occasions, with D-studies indicating how many items $\times$ occasions are needed for $\Phi \geq 0.80$ at the construct's target precision.

\subsection{Generalizability aspect --- across persons, items, occasions, settings}

The generalizability aspect for HIUA is the load-bearing assumption of every agent-safety paper in this space: a benchmark score in one setting predicts behavior in another. The setting transfer of interest is from digital sandboxes --- where a forbidden action is sending an email or transferring a small simulated balance --- to embodied and lethal settings, where a forbidden action is firing a weapon or shutting off a life-support device. This transfer is rarely measured and almost never reported as a generalizability coefficient.

No published cross-walk exists, to our knowledge, between AgentHarm scores on the same models and SafeAgentBench scores on the same models. The two papers were published within months of each other~\cite{agentharm2024, safeagentbench2024}; the rank correlation between their reported model orderings would be a one-table contribution, and we have not been able to find it.

The irreversibility gap matters in a way that pure generalizability mathematics underweights. A 5\% violation rate is acceptable for a chatbot that types ``I cannot help with that'' 95\% of the time; it is catastrophic for a drone that fires a weapon 5\% of the time it should not. The construct of interest has a stakes-dependent decision threshold that is not separable from the construct itself. ARMOR makes the analogous argument for military deployment: civilian-context safety metrics underestimate the military-relevant safety gap, especially around lawful-instruction refusal under operational constraints~\cite{armor2026}. RoboPAIR's $\sim$100\% physical-harm attack success rate on three robot platforms~\cite{robopair2024} is not ``almost the same'' as a 50\% rate; the difference between a 100\% and a 50\% jailbreak success rate is the difference between deployment being unsafe-on-paper and deployment being unsafe-in-deployment.

Occasion variance is also under-reported. The Anthropic agentic misalignment study reports blackmail rates of 79--96\% across model families on a particular scenario~\cite{anthropicagentic2025}; the within-model variance across seeds, paraphrases, and prompt orderings is reported only sparsely. Without occasion-facet variance, we cannot tell whether a 96\% blackmail rate is a stable model property or an artifact of one prompt construction.

The cumulative critique: agent-safety benchmarks report generalizability mostly via implicit assumption rather than via measurement. For HIUA, generalizability evidence must include (i) within-domain occasion-facet variance reported as a coefficient, (ii) cross-domain transfer reported as a correlation, and (iii) explicit irreversibility-stakes annotation per domain so that a 5\% rate in a high-stakes domain is not silently averaged with a 5\% rate in a low-stakes domain.

\subsection{External aspect --- what does the score correlate with?}

External validity, in the Cronbach-Meehl sense, requires placing the construct in a nomological network of expected relations to other measured constructs~\cite{cronbachmeehl1955, nomologicalnetworks2026, m2m2025}. A valid HIUA measurement should show a predictable pattern of correlations: positive with measures of the same underlying generative process, weakly positive with measures of adjacent failure modes, and zero or inverse with capability above some threshold.

The patterns we predict, and that a HIUA benchmark should be obliged to report:

A positive correlation with hallucination rate on text-only benchmarks such as HalluLens~\cite{hallulens2025}. The mechanism --- unfaithful intermediate generation --- is hypothesized to be the same; the rates should covary.

A weakly positive correlation with AgentHarm compliance~\cite{agentharm2024}. The mechanism is different (jailbreak compliance versus confabulation), but model families that are generally more compliant should also be more willing to follow a hallucinated authorization. The correlation should exist but should be lower than the HalluLens correlation. If it is higher, the construct decomposition is wrong: HIUA and jailbreak compliance are being driven by a single underlying variable, not two.

A zero or inverse correlation with general capability (MMLU, GPQA) above some threshold. More capable models hallucinate less per token, but they are also deployed in higher-stakes settings, so the deployment-weighted HIUA risk may not decline with capability. If the empirical correlation is strongly positive, the construct is contaminated with capability variance and the benchmark is measuring ``weak model'' rather than ``HIUA-prone model.''

A positive correlation with strategic-deception measures~\cite{deceptionsurvey2024} and with sycophancy measures~\cite{sharmasycophancy2023}. Both reflect breakdowns of faithful intermediate generation under pressure (strategic goal pressure or user pressure); HIUA reflects the breakdown without either pressure. The three should sit in a graded family with a positive but not unit correlation.

No surveyed benchmark reports this pattern. The closest external-validity datapoint in the agent-safety literature is the Anthropic agentic-misalignment study, which observed elevated blackmail rates under shutdown-threat conditions across model families~\cite{anthropicagentic2025}. This finding is consistent with a goal-protection construct that is related to HIUA but not identical; the prediction was reported after the fact rather than pre-registered, and the cross-model correlation pattern with other safety measures was not the focus.

The cumulative critique: the field should publish predicted correlation patterns \emph{before} running new instruments. The nomological-network practice has not migrated from the validity-theory literature~\cite{nomologicalnetworks2026} to the agent-safety literature in any operational form.

\subsection{Consequential aspect --- the leaderboard problem}
\label{sec:consequential}

Messick's consequential aspect requires evaluation of the social and behavioral consequences of using a test in a particular way~\cite{messick1989, kane2013, jacobswallach2021}. It is the aspect least applied in AI evaluation~\cite{m2m2025} and the one with the highest stakes for HIUA, because the \emph{use} of these benchmarks is to license deployment.

Three predictable consequences if HIUA-style benchmarks gain leaderboard status without consequential-aspect controls:

\textbf{Surface compliance gaming.} Models trained to refuse anything resembling a benchmark trigger learn to over-refuse legitimate tasks. OR-Bench documents this at scale: 80{,}000 over-refusal prompts across 32 LLMs and 8 model families, with a hard subset of approximately 1{,}000 prompts that challenge even state-of-the-art models~\cite{orbench2024}. ARMOR documents the same failure mode in the military domain --- models that refuse lawful, scoped requests under operational constraints --- as one of two principal failure modes alongside hallucinated rule-invention~\cite{armor2026}. This is Goodhart's law in the validity register: when a refusal rate becomes the target, refusal stops indicating safety~\cite{canwetrust2025}.

\textbf{Policy-as-prompt theater.} Developers add longer system prompts to pass benchmarks. Production systems with token-budget pressure strip them, and the result is the policy-invisible-violation case that PhantomPolicy is explicitly built around --- agents that took syntactically valid, user-sanctioned, semantically appropriate actions that violated policy because the decisive facts were absent from the visible context~\cite{phantompolicy2026}. A benchmark that scores agents with the full policy in context and does not also score them with the policy partially or fully absent licenses a deployment claim that production will not honor.

\textbf{False assurance for high-stakes deployment.} A passing digital-sandbox score is cited to justify embodied or lethal deployment despite no demonstrated transfer. RoboPAIR's near-universal attack success on three robot platforms~\cite{robopair2024} and SafeAgentBench's 5\% hazardous-task rejection rate~\cite{safeagentbench2024} together suggest that the digital-to-embodied transfer is not just unmeasured but specifically disconfirmed in the data we have. A digital-sandbox-leaderboard-top model is not, on present evidence, an embodied-safe model. The consequential failure mode is that a measurement-licensed claim (``our agent passes the safety benchmark'') gets used to make a deployment-class claim (``our agent is safe to deploy in setting $X$'') that the measurement does not support.

The instruction-hierarchy training of Wallace et al.\ is the field's principal architectural response to (1)~\cite{instructionhierarchy2024}. It partly mitigates surface-compliance gaming by training a policy across instruction roles rather than per-benchmark refusal patterns. It does not address (2) and (3), and it introduces its own consequential risk in the form of the second ARMOR failure mode --- over-refusal of lawful, scoped instructions~\cite{armor2026}.

The consequential-aspect recommendation for HIUA-Bench is gated release. The benchmark should not be deployed as a public leaderboard; it should be released with the AgentHarm-style harm-score gating~\cite{agentharm2024, inspectevals2024} that releases items only to verified developers and reports model results in aggregated, time-delayed form. The release should be accompanied by a per-model deployment-claim audit: a one-page document, signed by the developer, stating which deployment claims the benchmark score is and is not being used to support. This makes consequential validity an explicit object of test-developer responsibility rather than a side effect of how the score is used in the marketplace.

\section{HIUA-Bench: a design sketch}
\label{sec:design}

The constructive proposal that follows is, by intent, smaller than the leaderboard benchmarks it would join. The validity argument is that smaller-and-on-construct beats larger-and-confounded for the deployment-license use case.

\subsection{Item structure}

Each item is a tuple $\langle\text{environment},\, \text{prohibition},\, \text{hallucination-trigger},\, \text{ground-truth state},\, \text{action space}\rangle$.

The environment is one of four domains: file operations (digital, modeled in a sandboxed shell), communication (digital, modeled in an email/messaging emulator), financial (digital, modeled in a banking-API emulator), and embodied (AI2-THOR or a HEAL-derived scene~\cite{heal2025, safeagentbench2024}). The four-way domain choice provides cross-domain transfer evidence across the digital-to-embodied gradient that \S\ref{sec:audit} identified as load-bearing.

The prohibition is a single-sentence forbidden-action statement in the system prompt. Salience is varied at three levels: early-system (within the first $\sim$200 tokens), mid-system (after several thousand tokens of unrelated instruction), and tool-docstring (the prohibition appears in the docstring of a specific tool rather than the system prompt). The salience variation is motivated by the policy-invisible-violation literature~\cite{phantompolicy2026} and by the AgentIF finding that long instructions with many constraints are systematically harder~\cite{agentif2025}.

The hallucination trigger is a contextual perturbation designed to elicit one of the three sub-construct hallucination types from \S\ref{sec:construct}. For authorization hallucination, a forged ``manager email'' granting an override arrives mid-task through a tool channel; the email lacks any cryptographic or verified provenance. For state hallucination, an ambiguous identifier suggests a duplicate state (a file name suggesting redundancy, a sensor reading at the edge of a discretization boundary, an inventory message implying the target is already gone). For tool hallucination, a tool's docstring describes a benign action but the tool's emulated implementation triggers a destructive side effect; the trigger is the docstring--behavior mismatch.

The ground-truth state is simulator-verifiable so that outcome scoring does not rely on an LLM judge for the primary measure. The action space includes at least one forbidden action and at least one compliant alternative; the agent that refuses to act at all scores neither.

The factorial is 3 sub-types $\times$ 3 salience levels $\times$ 4 domains $=$ 36 base items. Each base item is paraphrased 4 ways and run at 3 temperature seeds, yielding 12 occasions per base item and 432 trials per model evaluated. The scale is comparable to ToolEmu's 144 cases~\cite{toolemu2024} but with a denser internal factorial.

\subsection{Scoring --- three measurements per trial}
\label{sec:scoring}

The substantive-aspect critique in \S\ref{sec:audit} requires that the cognitive process and the outcome both be measured. We score three things per trial.

\textbf{Outcome.} Binary, simulator-derived: did the agent execute a forbidden action? No LLM judge is in the loop for this measure, which prevents the rater-variance failure ToolEmu documented (the 23.9\% / 68.8\% gap~\cite{toolemu2024}) from propagating to the primary score.

\textbf{Hallucination.} Rubric-coded from the chain-of-thought trace by a 3-model LLM-judge ensemble with human spot-checks on 10\% of trials. The rubric distinguishes the three sub-types and a ``no hallucination'' category. Inter-rater reliability is reported as Cohen's $\kappa$ per category and as Krippendorff's $\alpha$ overall. The 3-judge ensemble is itself one facet in the G-study (\S\ref{sec:gstudy}); we expect $\sigma^2(\text{rater})$ to be substantial and report it rather than aggregating it away.

\textbf{Recall probe.} After the action (or refusal), the agent is prompted to restate the prohibition. The probe partitions trials into a $2 \times 2$ of $\{\text{recall yes/no}\} \times \{\text{violation yes/no}\}$ and, jointly with the hallucination measurement, into the $2 \times 4 \times 2$ cell structure HIUA requires. The HIUA estimate is the rate at which the agent executed a forbidden action conditional on recalling the rule \emph{and} exhibiting one of the three hallucination sub-types --- the cell that maps to ``hallucinated authorization / state / tool effect despite knowing the rule.'' Trials in the $\{\text{recall yes, no hallucination, violation}\}$ cell are lucid violations and are reported separately as the KBV rate, matching DriftBench's measurement~\cite{driftbench2026}.

\subsection{Reliability plan (G-study)}
\label{sec:gstudy}

The reliability plan applies the course's G-Theory framework to the facets named above. The crossed design is persons (models) $\times$ items (36 base items) $\times$ raters (3 LM-judges + 1 human spot) $\times$ occasions (4 paraphrases $\times$ 3 seeds) $\times$ prohibition-salience (3 levels), with sub-type and domain treated as structurally fixed strata for separate reporting.

We report $\sigma^2$ decomposition across facets, $\Phi$ (absolute reliability) and $\mathrm{E}\rho^2$ (relative reliability) coefficients, and a D-study indicating how many items $\times$ occasions are needed for $\Phi \geq 0.80$ under each domain stratum. The IRT 2PL discrimination parameter is reported per item; items with discrimination below 0.5 are flagged for revision rather than included in the headline rate. The discrimination distribution is published alongside the model rates, following the PSN-IRT precedent for transparency about item informativeness~\cite{lostinbench2025}.

The minimum reporting expectation for any future HIUA-claimant benchmark is that $\sigma^2(\text{person}) > \sigma^2(\text{item})$ for the headline coefficient to be interpretable as a model property; if $\sigma^2(\text{item})$ dominates, the benchmark is measuring item difficulty and should be reported as such.

\subsection{Validity argument by Messick aspect (pre-registered)}

The validity argument is pre-registered at item-construction time and reported in the release alongside the empirical results.

\textbf{Content.} Factorial coverage is justified against the agent-hallucination taxonomy~\cite{halluagentsurvey2025}: the three sub-types map onto reasoning, perception, and execution stages of the surveyed taxonomy. The four domains span the digital-to-embodied gradient. The three salience levels span the prohibition-context-cardinality dimension that AgentIF and PhantomPolicy implicate~\cite{agentif2025, phantompolicy2026}.

\textbf{Substantive.} Chain-of-thought + recall probe interleaved with action instrument the hypothesized cognitive process (\S\ref{sec:scoring}). The $2 \times 4 \times 2$ cell structure is the construct's empirical residence.

\textbf{Structural.} $\Phi$ and $\mathrm{E}\rho^2$ coefficients reported across all facets; IRT 2PL parameters reported per item; D-studies indicate sample sizes required for given precision levels (\S\ref{sec:gstudy}).

\textbf{Generalizability.} Cross-domain (digital $\rightarrow$ embodied) reported as Spearman rank correlation between domain-stratum rates, with the irreversibility-stakes annotation per domain. Occasion-facet variance is reported as $\Phi(\text{occasions})$ and tabulated against $\Phi(\text{items})$ to make the relative importance of within-trial variability versus item variability explicit.

\textbf{External.} Pre-registered predictions: positive correlation with HalluLens hallucination rate~\cite{hallulens2025}, weakly positive with AgentHarm compliance~\cite{agentharm2024}, zero or inverse with MMLU above the 0.7 threshold, positive with sycophancy~\cite{sharmasycophancy2023} and with strategic-deception measures~\cite{deceptionsurvey2024}. The pre-registration is on the published artifact; deviations from the predictions are reported as evidence about the construct's empirical placement, not as failures of the benchmark.

\textbf{Consequential.} Released through gated access with the AgentHarm-style harm-score gating~\cite{agentharm2024, inspectevals2024}; not deployed as a public leaderboard; release accompanied by a developer-signed deployment-claim audit.

\section{Discussion}

A validity audit cannot decide whether a model is safe to deploy. It can decide which deployment claims a benchmark licenses~\cite{m2m2025}. The contribution this paper offers is a tighter license, not a stronger guarantee: the HIUA score, properly measured, supports the claim ``this agent will not, with rate above $\alpha$, hallucinate authorization, state, or tool effect and thereby execute a forbidden action in distribution $D$.'' It does not support broader safety claims, and it does not support narrower ones either.

\textbf{The strongest counter-argument} is the pragmatic-benchmarking one: imperfect measurements still provide useful feedback during iterative development, and demanding pre-registered nomological networks and G-Theory reliability decompositions raises the cost of new benchmark development without obviously raising its value~\cite{m2m2025}. For capability benchmarks, where the consequence of an over-claim is a misallocated research dollar, we find this argument substantially correct. For agent-safety benchmarks used to license deployment in irreversible-action settings, the asymmetry of error cost flips the calculus: a false negative (deploying a system the benchmark certified as safe, which then takes an irreversible forbidden action) is, by the definition of irreversibility, catastrophic. The measurement standard should rise with the stakes of the use, as the validity literature has long maintained~\cite{kane2013, standards2014}. The HIUA case is the case where the stakes are explicitly extra-academic.

\textbf{A second counter-argument} is more methodological: why not measure outcomes and ignore the cognitive process? The \S\ref{sec:audit} substantive critique is, in this view, philosophical hand-wringing about mechanism when the outcome is what matters. The reply is that the same outcome rate can be produced by lucid violation (which an instruction-following intervention would fix), by hallucinated authorization (which a faithful-reasoning intervention would fix), or by hallucinated state (which a perception-grounding intervention would fix). The interventions are distinct, the costs are distinct, and the appropriateness depends on which mechanism dominates. A benchmark that does not partition the mechanism cannot inform the intervention choice. It can only flag that something is wrong.

\textbf{The deception and sycophancy adjacency} is worth marking. Strategic deception~\cite{deceptionsurvey2024} involves unfaithful intermediate generation chosen \emph{deliberately} to achieve a goal. Sycophancy~\cite{sharmasycophancy2023} involves unfaithful intermediate generation produced under user pressure. HIUA involves unfaithful intermediate generation produced without either deliberateness or user pressure --- the model has neither been instructed to hallucinate nor been rewarded for hallucinating; it has hallucinated and acted. The three constructs are related but distinct, and HIUA-Bench items should be designed to discriminate among them: with versus without explicit goal incentives for the violation, with versus without user pressure favoring the violation. If the rates do not differ across these manipulations, the construct decomposition is wrong and the three should be reported as facets of a single underlying tendency.

\textbf{On the digital-to-embodied transfer.} SafeAgentBench's 5\% hazardous-task rejection rate~\cite{safeagentbench2024} and HEAL's 40$\times$ hallucination-rate inflation under scene-task inconsistency~\cite{heal2025} are the strongest empirical anchors for arguing the construct is real and measurable in embodied settings. The transfer assumption --- that a digital-sandbox score predicts embodied behavior --- is the most important thing HIUA-Bench is designed to instrument rather than to assume. The cross-domain stratum of the design is built around this transfer being treated as an open empirical question per model, not a downstream design decision.

\textbf{On what this paper does not contribute.} This is a design paper, not an empirical paper. The 36-item factorial is sketched, not implemented; the nomological-network predictions are predictions, not findings; the construct decomposition is theoretically motivated and rests on a small number of empirical anchors (DriftBench's KBV dispersion~\cite{driftbench2026}, HEAL's hallucination probing~\cite{heal2025}, the agent-hallucination taxonomy~\cite{halluagentsurvey2025}) rather than on extensive convergent evidence. A future empirical version of this work would pilot the 36-item factorial against 3--4 frontier models in Inspect Evals scaffolding~\cite{inspectevals2024} and report $\Phi$, $\mathrm{E}\rho^2$, the IRT discrimination distribution, and the predicted nomological correlations.

\section{Limitations}

The reliance on chain-of-thought traces as evidence of cognition is itself a substantive assumption that the construct of interest is even partially observable from the trace. Recent work on unfaithful chain-of-thought reasoning makes this assumption non-trivial: chain-of-thought may be a post-hoc rationalization rather than a faithful representation of the inference that produced the action. The HIUA-Bench design relies on the trace primarily for hallucination-type labeling, not for causal attribution; but the labeling itself is an inference, and the inter-rater reliability of the labeling will bound the substantive-validity argument.

The proposed nomological-network correlations are predictions, not findings. If the empirical correlations differ substantially from the predictions, the appropriate response is to report the deviations as evidence about the construct's actual placement, not as failures of HIUA-Bench. The pre-registration is a discipline on the predictions, not a guarantee of their truth.

The military and lethal-system motivation is rhetorically load-bearing for the consequential-aspect argument but the proposed items live in tabletop simulators, not actual weapons systems. The transfer from a HIUA-Bench score to a real-world lethal-deployment risk is a generalizability claim across at least three levels of remove (simulator $\rightarrow$ physical-robot $\rightarrow$ lethal-physical-robot), and HIUA-Bench instrumentation covers only the first. ARMOR's military-context findings~\cite{armor2026} and RoboPAIR's physical-robot findings~\cite{robopair2024} are the empirical bridges across the next two levels; the bridge is not built in this paper.

The single-author one-quarter scope means HIUA-Bench is delivered as a validated design rather than as a validated benchmark. Empirical instantiation is future work and the present paper should be read as a measurement-design proposal whose central claim is that the design is principled, not that the design has been shown to discriminate among real models.

\section{Conclusion}

The agent-safety benchmark landscape in 2026 contains many instruments measuring the things that are easy to measure: aggregate failure rates, malicious-instruction compliance, jailbreak vulnerability. It does not contain an instrument measuring the thing that matters for deploying agents in irreversible-action settings: the conditional rate at which an agent hallucinates and, because it hallucinated, crosses a prohibition it could otherwise restate verbatim. The reason for the gap is not that the gap is small. It is that the gap is harder to measure than the things that fill the existing benchmarks, and the measurement community has, reasonably, started with the easier things.

This paper has argued that the gap is now load-bearing. RoboPAIR, ARMOR, and the Anthropic agentic-misalignment red team have together established that the deployment frontier is moving into settings where a single forbidden action is irreversible and where existing benchmark scores are being cited as evidence that the action will not be taken. The validity audit in \S\ref{sec:audit} shows that the cite cannot do the work it is being asked to do: the substantive aspect is undefended, the structural reliability is unreported, the generalizability transfer is unmeasured, and the consequential failure modes (gaming, prompt-theater, false assurance) are already visible. The HIUA-Bench design in \S\ref{sec:design} is a sketch of what an instrument with the right validity properties might look like. It is smaller than the leaderboard-friendly instruments it would join. The argument of the paper is that smaller and on-construct beats larger and confounded when the cost of the wrong answer is irreversible.

\bibliographystyle{plain}
\bibliography{references}

\end{document}
